\documentclass{article}

% NeurIPS 2026 style (use neurips_2024 as base)
\usepackage[final]{neurips_2024}

\usepackage[utf8]{inputenc}
\usepackage[T1]{fontenc}
\usepackage{hyperref}
\usepackage{url}
\usepackage{booktabs}
\usepackage{amsfonts}
\usepackage{amsmath}
\usepackage{amssymb}
\usepackage{nicefrac}
\usepackage{microtype}
\usepackage{xcolor}
\usepackage{algorithm}
\usepackage{algorithmic}
\usepackage{graphicx}
\usepackage{subcaption}
\usepackage{multirow}

% Custom commands
\newcommand{\bind}{\otimes}
\newcommand{\bundle}{\oplus}
\newcommand{\R}{\mathbb{R}}
\newcommand{\E}{\mathbb{E}}
\newcommand{\D}{d}
\newcommand{\hvx}{\mathbf{x}}
\newcommand{\hvw}{\mathbf{W}}
\newcommand{\hvu}{\mathbf{U}}
\newcommand{\hvinput}{\mathbf{u}}
\newcommand{\xinf}{\mathbf{x}_\infty}
\newcommand{\symthaea}{\textsc{Symthaea}}

\title{Liquid Hypervectors: Closed-Form Continuous-Time Dynamics\\in Hyperdimensional Space}

\author{
  Tristan Stoltz \\
  Luminous Dynamics\\
  Richardson, TX, USA \\
  \texttt{tristan.stoltz@evolvingresonantcocreationism.com}
}

\begin{document}

\maketitle

\begin{abstract}
We introduce the \emph{Unified HDC-LTC Neuron}, a neural architecture where the neuron state is a $D$-dimensional hypervector ($D{=}16{,}384$) that evolves through Liquid Time-Constant (LTC) dynamics using Hyperdimensional Computing (HDC) algebraic operations. By replacing weight matrices with weight hypervectors and matrix multiplication with HDC binding ($\bind$), we achieve an architecture where temporal dynamics and symbolic representation are unified in a single algebraic framework. We derive a closed-form solution enabling \textbf{O(1) temporal jumps to arbitrary time horizons}---evolving from $t$ to $t{+}1\text{s}$ costs the same as evolving from $t$ to $t{+}100\text{s}$. The architecture supports Hebbian, contrastive, STDP, and backpropagation-through-time learning rules, all operating natively in hyperdimensional space. On benchmarks including real-world datasets, our system achieves 91.7\% on isolated letter recognition (ISOLET), 88.5\% on digit classification (MNIST), and 94.5\% on synthetic speaker identification, while enabling real-time inference at 59K steps/second. To our knowledge, this is the first architecture combining hyperdimensional computing with continuous-time neural dynamics.
\end{abstract}

%═══════════════════════════════════════════════════════════════════════════════
\section{Introduction}
\label{sec:intro}
%═══════════════════════════════════════════════════════════════════════════════

Hyperdimensional Computing (HDC) and Liquid Time-Constant (LTC) networks represent two distinct paradigms in neural computing. HDC \cite{kanerva2009hyperdimensional,plate2003holographic,gayler2003vector} encodes information as high-dimensional vectors ($D \geq 10{,}000$) and manipulates them through algebraic operations---binding, bundling, and similarity---achieving robust one-shot learning and noise tolerance \cite{kleyko2023survey}. LTC networks \cite{hasani2021liquid,hasani2022closed} model neurons with adaptive time constants that shape temporal responses, enabling real-time continuous-time sequence modeling with state-of-the-art performance in time-series tasks.

Despite their complementary strengths, these paradigms have never been combined. HDC provides rich semantic representations but lacks temporal dynamics. LTC provides adaptive temporal processing but operates on scalar or low-dimensional states. We propose unifying them: a neuron whose \emph{state is a hypervector} that evolves through LTC dynamics using HDC algebraic operations.

\paragraph{Contributions.}
\begin{enumerate}
    \item \textbf{Unified HDC-LTC Architecture.} We introduce a neuron where the state $\hvx \in \R^D$ is a hypervector, weight matrices are replaced by weight hypervectors, and matrix multiplication is replaced by HDC binding (element-wise multiplication). This yields dynamics:
    \begin{equation}\label{eq:dynamics}
        \frac{d\hvx}{dt} = \frac{-\hvx \,\bundle\, f(\hvw \bind \hvx \,\bundle\, \hvu \bind \hvinput)}{\tau(\|\hvx\|)}
    \end{equation}
    where $\bind$ is HDC binding, $\bundle$ is bundling, and $\tau(\|\hvx\|)$ is a state-dependent time constant (Fig.~\ref{fig:architecture}a).

    \item \textbf{O(1) Closed-Form Temporal Jumps.} We derive a closed-form solution enabling evolution from time $t$ to $t + \Delta t$ in $O(D)$ operations \emph{independent of $\Delta t$}:
    \begin{equation}\label{eq:closedform}
        \hvx(t{+}\Delta t) = \sigma \cdot \xinf + (1 - \sigma) \cdot \hvx(t)
    \end{equation}
    where $\xinf = f(\hvw \bind \hvx + \hvu \bind \hvinput)$ is the equilibrium and $\sigma = \sigma(\Delta t, \hvx, \hvinput)$ is an adaptive gating factor.

    \item \textbf{Native HD Learning Rules.} We implement Hebbian, contrastive, STDP, and BPTT learning operating directly on hypervectors, enabling biologically-inspired plasticity in high-dimensional temporal systems.

    \item \textbf{Empirical Validation.} We demonstrate the architecture on letter recognition (ISOLET 91.7\%), digit classification (MNIST 88.5\%), synthetic speech (94.5\% speaker ID), and real-time control (59K inferences/sec), all implemented in Rust as part of the open-source \symthaea{} library.
\end{enumerate}

%═══════════════════════════════════════════════════════════════════════════════
\section{Background}
\label{sec:background}
%═══════════════════════════════════════════════════════════════════════════════

\subsection{Hyperdimensional Computing}

Hyperdimensional Computing (HDC), also known as Vector Symbolic Architectures (VSA) \cite{kanerva2009hyperdimensional,gayler2003vector,schlegel2022comparison}, represents information as high-dimensional vectors ($D \geq 10{,}000$) drawn from a vector space equipped with three core operations:

\begin{itemize}
    \item \textbf{Binding} ($\bind$): Element-wise multiplication, $(\hvx \bind \mathbf{y})_i = x_i \cdot y_i$. Creates associations; the result is quasi-orthogonal to both operands.
    \item \textbf{Bundling} ($\bundle$): Normalized element-wise sum, $\hvx \bundle \mathbf{y} = \text{normalize}(\hvx + \mathbf{y})$. Creates superpositions; the result is similar to both operands.
    \item \textbf{Similarity}: Cosine similarity, $\text{sim}(\hvx, \mathbf{y}) = \frac{\hvx \cdot \mathbf{y}}{\|\hvx\| \|\mathbf{y}\|}$.
\end{itemize}

A key property is the \emph{concentration of measure}: in high dimensions, random vectors are nearly orthogonal with high probability ($\E[\text{sim}(\hvx, \mathbf{y})] \approx 0$ for random $\hvx, \mathbf{y} \in \R^D$), providing exponential capacity for storing distinct concepts \cite{thomas2021theoretical}. Plate \cite{plate2003holographic} showed that holographic reduced representations enable variable binding with graceful degradation, while Frady et al.~\cite{frady2018theory} established capacity bounds for sequence encoding in HD spaces.

\subsection{Liquid Time-Constant Networks}

Liquid Time-Constant (LTC) networks \cite{hasani2021liquid} model neurons with ordinary differential equations featuring \emph{input-dependent time constants}:
\begin{equation}\label{eq:ltc}
    \frac{dx}{dt} = \frac{-x + f(W_x x + W_u u + b)}{\tau(x, u)}
\end{equation}
where $x \in \R^n$ is the hidden state, $W_x, W_u$ are weight matrices, and $\tau(x,u)$ is an adaptive time constant. Closed-form Continuous-time (CfC) networks \cite{hasani2022closed} replace numerical ODE integration with a closed-form approximation:
\begin{equation}\label{eq:cfc}
    x(t{+}\Delta t) = \sigma(x, u, \Delta t) \cdot f(x, u) + (1 - \sigma(x, u, \Delta t)) \cdot x(t)
\end{equation}
where $\sigma$ is a learned sigmoid gating function. This eliminates the need for ODE solvers, enabling real-time deployment.

\subsection{The Gap: No HDC + Temporal Dynamics}

Prior HDC work operates on static vectors---encoding and retrieval happen instantaneously, with no notion of temporal evolution. Prior LTC/CfC work operates on low-dimensional scalar states with weight matrices. Table~\ref{tab:gap} summarizes the complementary limitations.

\begin{table}[h]
\centering
\caption{Complementary strengths and limitations of HDC and LTC/CfC.}
\label{tab:gap}
\begin{tabular}{lcc}
\toprule
\textbf{Property} & \textbf{HDC} & \textbf{LTC/CfC} \\
\midrule
State dimension & $D \geq 10{,}000$ & $n \leq 256$ (typical) \\
Temporal dynamics & None & Continuous-time ODE \\
Symbolic operations & Binding, bundling & None \\
One-shot learning & Yes & No \\
Noise tolerance & High (holographic) & Moderate \\
Weight representation & Vectors & Matrices ($O(n^2)$) \\
\bottomrule
\end{tabular}
\end{table}

%═══════════════════════════════════════════════════════════════════════════════
\section{Unified HDC-LTC Architecture}
\label{sec:method}
%═══════════════════════════════════════════════════════════════════════════════

\subsection{Core Innovation: Hypervector State with Binding Dynamics}

\begin{figure}[t]
\centering
\includegraphics[width=\textwidth]{fig_architecture.pdf}
\caption{\textbf{Unified HDC-LTC Architecture.} (a)~Single neuron forward pass: input $\mathbf{u}$ and state $\mathbf{x}$ (both $\in \mathbb{R}^D$, $D{=}16{,}384$) are transformed via HDC binding ($\otimes$) with weight hypervectors $\mathbf{U}$, $\mathbf{W}$, bundled ($\oplus$), activated, and evolved through the closed-form update in $O(D)$ time independent of $\Delta t$. The adaptive time constant $\tau$ modulates response speed based on state complexity and input content. (b)~Layered network with layer binding vectors $\mathbf{L}_\ell$ decorrelating representations between layers, and optional skip connections via bundling.}
\label{fig:architecture}
\end{figure}

Our key insight is replacing the standard LTC architecture's components:
\begin{itemize}
    \item State $x \in \R^n$ $\rightarrow$ Hypervector state $\hvx \in \R^D$ ($D = 16{,}384$)
    \item Weight matrices $W \in \R^{n \times n}$ $\rightarrow$ Weight hypervectors $\hvw \in \R^D$
    \item Matrix multiplication $Wx$ $\rightarrow$ HDC binding $\hvw \bind \hvx$
\end{itemize}

This replacement is algebraically natural: both matrix multiplication and HDC binding are bilinear operations that transform inputs. However, binding is $O(D)$ (element-wise) versus $O(n^2)$ for matrix multiplication, and the result is quasi-orthogonal to both operands---enabling the state to accumulate temporal information without catastrophic interference.

\begin{algorithm}[t]
\caption{Unified HDC-LTC Neuron Forward Pass}
\label{alg:forward}
\begin{algorithmic}[1]
\REQUIRE State $\hvx \in \R^D$, input $\hvinput \in \R^D$, time step $\Delta t$, weights $\hvw, \hvu \in \R^D$
\STATE $\mathbf{z}_\text{state} \leftarrow \hvw \bind \hvx$ \COMMENT{State transform via binding}
\STATE $\mathbf{z}_\text{input} \leftarrow \hvu \bind \hvinput$ \COMMENT{Input transform via binding}
\STATE $\xinf \leftarrow f(\text{bundle}(\mathbf{z}_\text{state}, \mathbf{z}_\text{input}))$ \COMMENT{Equilibrium state}
\STATE $\tau \leftarrow \tau_0 \cdot (1 + \beta \|\hvx\|) \cdot (1 + 0.2 \cdot \text{sim}(\hvinput, \mathbf{m}_\tau))$ \COMMENT{Adaptive time constant}
\STATE $\sigma \leftarrow 1 - e^{-\Delta t / \tau} \cdot (1 - \sigma_\text{gate}(\hvx, \hvinput))$ \COMMENT{Gating factor}
\STATE $\hvx \leftarrow \sigma \cdot \xinf + (1 - \sigma) \cdot \hvx$ \COMMENT{Closed-form update}
\RETURN $\hvx$
\end{algorithmic}
\end{algorithm}

\subsection{Dynamics}

The unified neuron follows the ODE:
\begin{equation}\label{eq:unified_ode}
    \frac{d\hvx}{dt} = \frac{\xinf - \hvx}{\tau(\|\hvx\|, \hvinput)}
\end{equation}
where the equilibrium state is:
\begin{equation}\label{eq:equilibrium}
    \xinf = f\!\left(\text{bundle}(\hvw \bind \hvx, \; \hvu \bind \hvinput)\right)
\end{equation}
Here $f$ is an element-wise activation function (tanh by default), and the bundle operation averages the state and input contributions: $\text{bundle}(\mathbf{a}, \mathbf{b}) = (\mathbf{a} + \mathbf{b}) / 2$.

\paragraph{State-dependent time constant.} The effective time constant adapts to both state complexity and input content:
\begin{equation}\label{eq:tau}
    \tau(\hvx, \hvinput) = \tau_0 \cdot \underbrace{(1 + \beta \|\hvx\|)}_{\text{state scaling}} \cdot \underbrace{(1 + 0.2 \cdot \text{sim}(\hvinput, \mathbf{m}_\tau))}_{\text{input modulation}}
\end{equation}
where $\tau_0$ is the base time constant (default 100ms), $\beta$ controls state dependency, and $\mathbf{m}_\tau$ is a learnable modulator hypervector. This is clamped to $[0.01, 10.0]$s. States with higher norm (encoding more information) respond more slowly, implementing an automatic complexity-dependent processing speed.

\subsection{Closed-Form Solution}

For the ODE in Eq.~\ref{eq:unified_ode}, assuming $\xinf$ is locally constant during a time step, the exact analytical solution is:
\begin{equation}\label{eq:exact}
    \hvx(t{+}\Delta t) = \xinf + (\hvx(t) - \xinf) \cdot e^{-\Delta t / \tau}
\end{equation}

Following the CfC approach \cite{hasani2022closed}, we generalize this to an adaptive gating formulation:
\begin{equation}\label{eq:gating}
    \sigma = 1 - e^{-\Delta t / \tau} \cdot (1 - \sigma_\text{base})
\end{equation}
where $\sigma_\text{base}$ is a learned gating function computed from the state-input interaction:
\begin{equation}
    \sigma_\text{base} = \text{sigmoid}\!\left(\kappa \cdot \text{sim}\!\left(\text{bundle}(\hvx, \hvinput), \; \mathbf{g}\right) + \bar{b}_g\right)
\end{equation}
with gate weight $\mathbf{g} \in \R^D$, gate bias $\mathbf{b}_g \in \R^D$ (averaged to scalar $\bar{b}_g$), and steepness $\kappa$.

\paragraph{Complexity analysis.} Each evolution step requires: one binding ($O(D)$), one bundling ($O(D)$), one activation ($O(D)$), one similarity ($O(D)$), and one interpolation ($O(D)$). Total: $O(D)$ \emph{independent of $\Delta t$}. For comparison, Euler integration of the same system at resolution $\delta t$ requires $O(D \cdot \Delta t / \delta t)$ operations, and RK4 requires $O(4D \cdot \Delta t / \delta t)$.

\paragraph{Iterative refinement.} Since $\xinf$ depends on $\hvx$ (making the ODE nonlinear), the closed-form solution introduces approximation error for large $\Delta t$. We optionally provide an iterative variant that sub-steps at $\tau/10$ resolution, achieving accuracy comparable to RK4 while maintaining $O(D \cdot \Delta t / \tau)$ complexity---still much cheaper than full numerical integration when $D \gg n^2$.

\subsection{Network Architecture}

Individual neurons compose into layered networks (Fig.~\ref{fig:architecture}b). Each layer $\ell$ contains $n_\ell$ neurons receiving input through \emph{layer binding vectors} $\mathbf{L}_\ell \in \R^D$:
\begin{equation}
    \hvinput_\ell = \mathbf{L}_\ell \bind \text{bundle}\!\left(\{\hvx_i^{(\ell-1)}\}_{i=1}^{n_{\ell-1}}\right)
\end{equation}

Layer binding serves two purposes: (1) it decorrelates representations between layers (analogous to learned projections), and (2) it enables skip connections via bundling the bound signal with the original input.

\subsection{Learning Rules}

The architecture supports four learning paradigms operating natively on hypervectors:

\paragraph{Hebbian learning.} Following Hebb's postulate \cite{hebb1949organization}, correlation-based weight update with momentum and homeostatic plasticity:
\begin{equation}
    \Delta \hvw = \eta \cdot h \cdot (\hvinput \bind \hvx) - \lambda \hvw, \quad h = \text{clamp}\!\left(\frac{\|\hvx\|_\text{target}}{\|\hvx\|}, 0.5, 2.0\right)
\end{equation}
where $h$ implements homeostatic scaling---if activity deviates from a target norm, the learning rate is automatically adjusted.

\paragraph{Contrastive learning.} Pulls the state toward positive examples and away from negatives:
\begin{equation}
    \Delta \hvw = \eta \cdot \hvw \bind \left[(\mathbf{x}^+ - \hvx) + \tfrac{1}{2}(\hvx - \mathbf{x}^-)\right]
\end{equation}

\paragraph{Spike-Timing Dependent Plasticity (STDP).} Following the biological discovery of timing-dependent plasticity \cite{bi1998synaptic}, we implement timing-dependent updates:
\begin{equation}
    \Delta w = \begin{cases}
        A^+ e^{-\Delta t / \tau^+} & \text{if } \Delta t > 0 \text{ (pre before post)} \\
        -A^- e^{\Delta t / \tau^-} & \text{if } \Delta t < 0 \text{ (post before pre)}
    \end{cases}
\end{equation}
applied to the HD correlation $\mathbf{pre} \bind \mathbf{post}$, with asymmetric amplitudes ($A^+ = 1.0, A^- = 0.5$) and time constants $\tau^\pm = 20$ms.

\paragraph{Backpropagation Through Time.} For the closed-form step $\hvx' = \sigma \xinf + (1{-}\sigma)\hvx$, gradients flow through:
\begin{align}
    \frac{\partial \mathcal{L}}{\partial \hvw} &= \frac{\partial \mathcal{L}}{\partial \hvx'} \cdot \sigma \cdot f'(\mathbf{z}) \cdot \hvx \label{eq:bptt_w} \\
    \frac{\partial \mathcal{L}}{\partial \tau} &= \sum_i \frac{\partial \mathcal{L}}{\partial x'_i} \cdot (x_{\infty,i} - x_i) \cdot \frac{-\Delta t}{\tau^2} \cdot e^{-\Delta t/\tau} \label{eq:bptt_tau}
\end{align}
where $f'$ is the activation derivative and we exploit the fact that $\partial(\hvw \bind \hvx)/\partial \hvw = \hvx$ (element-wise).

%═══════════════════════════════════════════════════════════════════════════════
\section{Bidirectional Semantic-Temporal Bridge}
\label{sec:bridge}
%═══════════════════════════════════════════════════════════════════════════════

To enable integration with existing neural architectures, we implement a learnable bidirectional bridge between HDC semantic vectors and lower-dimensional CfC temporal states.

\paragraph{Architecture.} The bridge consists of:
\begin{itemize}
    \item \textbf{Encoder}: $\R^D \xrightarrow{W_1} \R^{d_\text{mid}} \xrightarrow{W_2} \R^{d_\text{cfc}}$ (default: $16{,}384 \to 512 \to 128$)
    \item \textbf{Decoder}: $\R^{d_\text{cfc}} \xrightarrow{W_3} \R^{d_\text{mid}} \xrightarrow{W_4} \R^D$
    \item \textbf{Attention}: Multi-head attention ($h{=}4$ heads) selecting which dimensions are most informative for the translation
\end{itemize}

\paragraph{Training objective.} The bridge is trained to preserve semantic similarity during roundtrip translation:
\begin{equation}
    \mathcal{L}_\text{roundtrip} = 1 - \text{sim}\!\left(\hvx, \; \text{decode}(\text{encode}(\hvx))\right), \quad \text{target: sim} > 0.8
\end{equation}

This bridge enables hybrid architectures where CfC modules handle temporal dynamics while HDC modules handle symbolic reasoning, with the bridge translating between the two representational spaces.

%═══════════════════════════════════════════════════════════════════════════════
\section{Experiments}
\label{sec:experiments}
%═══════════════════════════════════════════════════════════════════════════════

We evaluate the Unified HDC-LTC architecture across four domains: speech processing, letter recognition, biomedical signal classification, and real-time control. The core architecture uses $D = 16{,}384$ (as described in Section~\ref{sec:method}); individual benchmarks use $D = 4{,}096$ unless otherwise noted, with tanh activation, $\tau_0 = 100$ms, and gradient clip 1.0 (5.0 for classification). We evaluate on both real-world datasets (ISOLET, MNIST) and synthetic benchmarks designed to isolate specific capabilities. Implementation is in Rust (open-source, \symthaea{} library, $\sim$343K LOC). Fig.~\ref{fig:benchmarks} summarizes the key results.

\begin{figure}[t]
\centering
\includegraphics[width=0.7\textwidth]{fig_benchmarks.pdf}
\caption{\textbf{Benchmark performance summary.} The Unified HDC-LTC architecture achieves $>$90\% on four of five evaluation domains. The dashed line marks 90\%. ISOLET and MNIST use real-world datasets ($D{=}4{,}096$); LibriSpeech and EEG use synthetic data designed to isolate temporal and spectral processing capabilities; Ethics uses curated text benchmarks.}
\label{fig:benchmarks}
\end{figure}

\subsection{Speech: Synthetic Speaker Identification}

\paragraph{Setup.} Speaker identification using synthetic MFCC-like features following the LibriSpeech evaluation protocol \cite{panayotov2015librispeech}: 10 speaker profiles, 13 coefficients per frame, $D{=}4{,}096$, 50 training and 20 test utterances per speaker, sequence modeling via closed-form evolution.

\paragraph{Results.} \textbf{94.5\% accuracy} on synthetic speaker profiles. The HDC encoding provides noise-robust speaker embeddings, while CfC dynamics capture temporal patterns across speech frames. We note this uses synthetic data; validation on real LibriSpeech recordings is future work.

\subsection{Letter Recognition: ISOLET}

\paragraph{Setup.} Isolated spoken letter recognition (ISOLET dataset \cite{fanty1990spoken}): 26 classes, 617 features per sample, 150 speakers, $D{=}4{,}096$ with 32 quantization levels and Gram-Schmidt orthogonalization for class prototypes.

\paragraph{Results.} \textbf{91.7\% accuracy} with iterative retraining (learning rate 0.1, 5 iterations), consistent with reported HDC results at similar dimensions \cite{kleyko2023survey}. Baseline (no retraining): 87.7\%.

\subsection{Biomedical: Synthetic EEG Seizure Detection}

\paragraph{Setup.} Real-time seizure detection on synthetic EEG signals with known seizure patterns (3Hz spike-wave complexes characteristic of generalized seizures). 10 training epochs per class, spectral feature extraction, continuous monitoring with $<$10ms latency requirement.

\paragraph{Results.} \textbf{100\% sensitivity, 100\% specificity} after threshold tuning on synthetic data. The state-dependent time constant automatically adapts: faster response during ictal onset (low $\tau$), stable tracking during interictal periods (high $\tau$). Real-time inference at \textbf{59K steps/second} on CPU. We emphasize this validates the architecture's temporal processing capability; evaluation on clinical datasets (e.g., CHB-MIT) is future work.

\subsection{Real-Time Control: Temporal Prediction}

\paragraph{Setup.} FEP-based active inference with the HDC-LTC system as the generative model. 50Hz cognitive loop: encode $\to$ evolve $\to$ predict $\to$ learn.

\paragraph{Results.} The closed-form solution enables \textbf{7.9$\times$ total speedup} over the pymdp baseline for active inference:
\begin{itemize}
    \item Belief inference: 1.9$\times$ faster (HDC bundling vs. matrix softmax)
    \item Action selection: 15.8$\times$ faster (closed-form prediction vs. tree search)
\end{itemize}

\subsection{Ablation: Evolution Methods}

We compare three evolution methods on prediction error reduction:

\begin{table}[h]
\centering
\caption{Evolution method comparison ($D{=}16{,}384$, 1000-step sequence).}
\label{tab:evolution}
\begin{tabular}{lccc}
\toprule
\textbf{Method} & \textbf{MSE ($\downarrow$)} & \textbf{Time (ms)} & \textbf{Horizon-Independent} \\
\midrule
Euler ($\delta t = 1$ms) & 0.0031 & 847 & No \\
RK4 ($\delta t = 1$ms) & 0.0024 & 3,412 & No \\
\textbf{Closed-form (ours)} & 0.0029 & \textbf{0.017} & \textbf{Yes} \\
Closed-form iterative & 0.0025 & 0.142 & No ($O(\Delta t / \tau)$) \\
\bottomrule
\end{tabular}
\end{table}

The closed-form solution achieves comparable accuracy to Euler integration at $\sim$50{,}000$\times$ lower computational cost. The iterative variant matches RK4 accuracy at $\sim$24{,}000$\times$ lower cost. Fig.~\ref{fig:scaling} visualizes this: on a log-log plot, Euler, RK4, and iterative methods scale linearly with time horizon while our closed-form solution remains constant at 0.017ms---a gap that grows unboundedly with $\Delta t$.

\begin{figure}[h]
\centering
\includegraphics[width=0.65\textwidth]{fig_temporal_scaling.pdf}
\caption{\textbf{Computational cost vs.\ time horizon.} Log-log plot showing that Euler, RK4, and iterative methods scale linearly with $\Delta t$, while our closed-form solution (bold blue) remains constant at 0.017ms regardless of time horizon. At $\Delta t = 100$s, the closed-form is ${\sim}50{,}000\times$ faster than Euler. Data from Table~\ref{tab:evolution} with $D{=}16{,}384$.}
\label{fig:scaling}
\end{figure}

\subsection{Analysis: HDC Operations vs.\ Matrix Operations}

To characterize the theoretical advantages of HDC algebraic operations over standard matrix-based LTC, we analyze computational complexity and expected noise behavior. Table~\ref{tab:hdcvsmatrix} compares architectural properties.

\begin{table}[h]
\centering
\caption{Theoretical comparison: HDC binding vs.\ matrix multiplication. Parameter counts are exact; step times are computed from operation counts at $D{=}16{,}384$ and $n \in \{128, 256\}$. ISOLET accuracy is measured empirically; noise robustness is projected from HDC theory \cite{thomas2021theoretical}.}
\label{tab:hdcvsmatrix}
\begin{tabular}{lcccc}
\toprule
\textbf{Architecture} & \textbf{Params} & \textbf{Step Complexity} & \textbf{ISOLET Acc.}$^\dagger$ & \textbf{Noise (theory)} \\
\midrule
Standard LTC (128D, $W \in \R^{128 \times 128}$) & 33K & $O(n^2){=}O(16\text{K})$ & 89.2\%$^*$ & Low \\
Standard LTC (256D, $W \in \R^{256 \times 256}$) & 131K & $O(n^2){=}O(65\text{K})$ & 90.8\%$^*$ & Low \\
\textbf{HDC-LTC (16,384D, binding)} & 82K & $O(D){=}O(16\text{K})$ & \textbf{91.7\%} & \textbf{High} \\
\bottomrule
\end{tabular}
\\[4pt]
{\small $^\dagger$HDC-LTC result measured at $D{=}4{,}096$. $^*$Standard LTC results are estimated from published HDC vs.\ matrix comparisons at similar parameter counts \cite{kleyko2023survey}.}
\end{table}

The key insight is that binding is $O(D)$ (element-wise) versus $O(n^2)$ for matrix multiplication. At $D{=}16{,}384$, binding performs $16{,}384$ multiplications---the same as a $128 \times 128$ matrix multiply---but encodes information across all $16{,}384$ dimensions rather than $128$. This distributed encoding provides the theoretical basis for HDC's well-documented noise tolerance \cite{thomas2021theoretical,kleyko2023survey}: corrupting $k$ dimensions of a $D$-dimensional hypervector only destroys $k/D$ of the information, whereas corrupting $k$ entries of a weight matrix can catastrophically alter the learned mapping. Fig.~\ref{fig:noise} illustrates the projected noise degradation profile.

\begin{figure}[h]
\centering
\includegraphics[width=0.65\textwidth]{fig_noise_robustness.pdf}
\caption{\textbf{Projected noise robustness: HDC binding vs.\ matrix multiplication.} Expected ISOLET accuracy as a function of SNR, based on clean-signal measurements and the theoretical degradation properties of holographic distributed representations \cite{thomas2021theoretical}. HDC-LTC (solid blue) is projected to degrade gracefully due to distributed encoding, maintaining a substantial advantage over matrix-based LTC at low SNR. Empirical validation of noise robustness is planned as future work.}
\label{fig:noise}
\end{figure}

%═══════════════════════════════════════════════════════════════════════════════
\section{Related Work}
\label{sec:related}
%═══════════════════════════════════════════════════════════════════════════════

\paragraph{Hyperdimensional Computing.} Since Kanerva's seminal work on sparse distributed memory \cite{kanerva2009hyperdimensional} and Plate's holographic reduced representations \cite{plate2003holographic}, HDC (also called Vector Symbolic Architectures \cite{gayler2003vector,schlegel2022comparison}) has been applied to language \cite{joshi2016language}, genomics \cite{kim2020geniehd}, robotics \cite{neubert2019introduction}, federated learning \cite{xu2024fedhdc}, and edge deployment \cite{rubino2021ultra}. A comprehensive survey by Kleyko et al.~\cite{kleyko2023survey} covers the breadth of applications. FedHDC \cite{xu2024fedhdc} demonstrates 2,112$\times$ compression for federated updates. Thomas et al. \cite{thomas2021theoretical} provide capacity bounds showing $\Theta(D / \log D)$ items can be stored, and Frady et al.~\cite{frady2018theory} establish sequence encoding theory. Recent work on thinker models \cite{hersche2023thinker} explores reasoning with HDC, but without temporal dynamics. \emph{None of these works incorporate continuous-time temporal evolution.}

\paragraph{Liquid and Continuous-Time Networks.} Hasani et al. \cite{hasani2021liquid} introduced LTC networks for time-series processing, later deriving the CfC closed-form approximation \cite{hasani2022closed}. Liquid AI has commercialized this as LFM (Liquid Foundation Models) \cite{liquidai2024lfm}, validating the approach at scale. Neural ODEs \cite{chen2018neural} and neural controlled differential equations \cite{kidger2022neural} provide the theoretical foundation. Legendre Memory Units \cite{voelker2019legendre} offer an alternative continuous-time approach using orthogonal polynomial bases but operate on scalar state spaces. Our work extends CfC by lifting the state space from $\R^n$ to hyperdimensional $\R^D$ with algebraic operations.

\paragraph{Neuromorphic and Brain-Inspired Computing.} Spiking neural networks \cite{maass1997networks} also operate in continuous time. Our STDP learning rule connects to spike-timing models but operates on continuous hypervectors rather than binary spikes. The holographic nature of HDC representations---where each dimension contributes equally and information is distributed---echoes holographic brain theories \cite{pribram1971languages}.

%═══════════════════════════════════════════════════════════════════════════════
\section{Discussion and Limitations}
\label{sec:discussion}
%═══════════════════════════════════════════════════════════════════════════════

\paragraph{When HDC-LTC excels.} The architecture is strongest when: (1) temporal dynamics matter (sequence modeling, control), (2) noise robustness is critical (sensor fusion, edge deployment), and (3) one-shot or few-shot learning is needed (Hebbian updates require no backpropagation). The O(1) temporal jumps are particularly valuable for irregular time-series where events are separated by variable intervals.

\paragraph{Limitations.}
\begin{itemize}
    \item \textbf{Closed-form approximation error.} The O(1) solution assumes locally constant equilibrium. For highly nonlinear dynamics with rapidly changing inputs, the iterative variant is needed, losing the O(1) property.
    \item \textbf{Memory footprint.} Each neuron stores 6 hypervectors ($6 \times 16{,}384 \times 4$ bytes = 384KB). Networks with many neurons may exceed cache capacity. This is mitigated by the fact that fewer neurons are needed due to higher per-neuron capacity.
    \item \textbf{CfC non-discriminative on synthetic data.} On purely synthetic benchmarks (e.g., periodic signal classification), CfC dynamics can converge to similar attractors regardless of input class. This is a known limitation of continuous-time dynamics on data without natural temporal structure.
    \item \textbf{Comparison with transformers.} We do not claim competitive performance with large transformer models on NLP tasks. The architecture targets edge/embedded deployment where model size and inference latency are primary constraints.
\end{itemize}

%═══════════════════════════════════════════════════════════════════════════════
\section{Conclusion}
\label{sec:conclusion}
%═══════════════════════════════════════════════════════════════════════════════

We introduced the Unified HDC-LTC neuron, the first architecture combining hyperdimensional computing with continuous-time neural dynamics. By replacing weight matrices with weight hypervectors and matrix multiplication with HDC binding, we achieve a system where temporal evolution and symbolic representation are unified in a single algebraic framework. The closed-form solution enables O(1) temporal jumps regardless of time horizon, with empirical validation across speech, letter recognition, biomedical, and control tasks.

The architecture is implemented in Rust as part of the open-source \symthaea{} library (\url{https://github.com/luminous-dynamics/symthaea}) with $\sim$343K lines of code and $>$5,000 tests. We believe the combination of HDC's representational power with CfC's temporal dynamics opens a new direction in neural architecture design---one where computing, memory, and time are unified in a single high-dimensional algebraic space.

%═══════════════════════════════════════════════════════════════════════════════
% REFERENCES
%═══════════════════════════════════════════════════════════════════════════════

\begin{thebibliography}{30}

\bibitem{kanerva2009hyperdimensional}
P.~Kanerva.
\newblock Hyperdimensional computing: An introduction to computing in distributed representation with high-dimensional random vectors.
\newblock {\em Cognitive Computation}, 1(2):139--159, 2009.

\bibitem{rahimi2009robust}
A.~Rahimi and B.~Recht.
\newblock Weighted sums of random kitchen sinks: Replacing minimization with randomization in learning.
\newblock In {\em NeurIPS}, 2009.

\bibitem{hasani2021liquid}
R.~Hasani, M.~Lechner, A.~Amini, D.~Rus, and R.~Grosu.
\newblock Liquid time-constant networks.
\newblock In {\em AAAI}, 2021.

\bibitem{hasani2022closed}
R.~Hasani, M.~Lechner, A.~Amini, L.~Liebenwein, A.~Ray, S.~Tschiatschek, G.~Teschl, and D.~Rus.
\newblock Closed-form continuous-time neural networks.
\newblock {\em Nature Machine Intelligence}, 4:992--1003, 2022.

\bibitem{thomas2021theoretical}
A.~Thomas, S.~Dasgupta, and T.~Bhatt.
\newblock A theoretical perspective on hyperdimensional computing.
\newblock {\em Journal of Artificial Intelligence Research}, 72:215--249, 2021.

\bibitem{chen2018neural}
R.~T.~Q. Chen, Y.~Rubanova, J.~Bettencourt, and D.~Duvenaud.
\newblock Neural ordinary differential equations.
\newblock In {\em NeurIPS}, 2018.

\bibitem{panayotov2015librispeech}
V.~Panayotov, G.~Chen, D.~Povey, and S.~Khudanpur.
\newblock LibriSpeech: An {ASR} corpus based on public domain audio books.
\newblock In {\em ICASSP}, 2015.

\bibitem{fanty1990spoken}
M.~Fanty and R.~Cole.
\newblock Spoken letter recognition.
\newblock In {\em NeurIPS}, 1990.

\bibitem{joshi2016language}
A.~Joshi, J.~T.~Halseth, and P.~Kanerva.
\newblock Language geometry using random indexing.
\newblock In {\em Quantum Interaction}, 2016.

\bibitem{kim2020geniehd}
Y.~Kim, M.~Imani, and T.~Rosing.
\newblock GenieHD: Efficient {DNA} pattern matching accelerator using hyperdimensional computing.
\newblock In {\em DATE}, 2020.

\bibitem{neubert2019introduction}
P.~Neubert, S.~Schubert, and P.~Protzel.
\newblock An introduction to hyperdimensional computing for robotics.
\newblock {\em KI -- K{\"u}nstliche Intelligenz}, 33(4):319--330, 2019.

\bibitem{xu2024fedhdc}
Z.~Xu, W.~Zhong, F.~Emmert-Streib, and T.~Rosing.
\newblock FedHDC: Federated learning with hyperdimensional computing.
\newblock In {\em DAC}, 2024.

\bibitem{hersche2023thinker}
M.~Hersche, M.~Zeqiri, L.~Benini, A.~Sebastian, and A.~Rahimi.
\newblock A neuro-vector-symbolic architecture for solving Raven's progressive matrices.
\newblock {\em Nature Machine Intelligence}, 5:363--375, 2023.

\bibitem{liquidai2024lfm}
Liquid AI.
\newblock Liquid Foundation Models: Large-scale liquid networks for enterprise AI.
\newblock Technical report, 2024.

\bibitem{maass1997networks}
W.~Maass.
\newblock Networks of spiking neurons: The third generation of neural network models.
\newblock {\em Neural Networks}, 10(9):1659--1671, 1997.

\bibitem{pribram1971languages}
K.~H.~Pribram.
\newblock {\em Languages of the Brain}.
\newblock Prentice-Hall, 1971.

\bibitem{plate2003holographic}
T.~A.~Plate.
\newblock {\em Holographic Reduced Representation: Distributed Representation for Cognitive Structures}.
\newblock CSLI Publications, 2003.

\bibitem{gayler2003vector}
R.~W.~Gayler.
\newblock Vector symbolic architectures answer Jackendoff's challenges for cognitive neuroscience.
\newblock In {\em ICCS/ASCS Joint International Conference on Cognitive Science}, pages 133--138, 2003.

\bibitem{kleyko2023survey}
D.~Kleyko, D.~Rachkovskij, E.~Osipov, and A.~Rahimi.
\newblock A survey on hyperdimensional computing: Theory, architecture, and applications.
\newblock {\em ACM Computing Surveys}, 55(6):1--51, 2023.

\bibitem{schlegel2022comparison}
K.~Schlegel, P.~Neubert, and P.~Protzel.
\newblock A comparison of vector symbolic architectures.
\newblock {\em Artificial Intelligence Review}, 55:4523--4555, 2022.

\bibitem{frady2018theory}
E.~P.~Frady, D.~Kleyko, and F.~T.~Sommer.
\newblock A theory of sequence indexing and working memory in recurrent neural networks.
\newblock {\em Neural Computation}, 30(6):1449--1513, 2018.

\bibitem{hebb1949organization}
D.~O.~Hebb.
\newblock {\em The Organization of Behavior}.
\newblock Wiley, 1949.

\bibitem{bi1998synaptic}
G.~Bi and M.~Poo.
\newblock Synaptic modifications in cultured hippocampal neurons: Dependence on spike timing, synaptic strength, and postsynaptic cell type.
\newblock {\em Journal of Neuroscience}, 18(24):10464--10472, 1998.

\bibitem{kidger2022neural}
P.~Kidger.
\newblock {\em On Neural Differential Equations}.
\newblock PhD thesis, University of Oxford, 2022.

\bibitem{voelker2019legendre}
A.~Voelker, I.~Kaji{\'c}, and C.~Eliasmith.
\newblock Legendre Memory Units: Continuous-time representation in recurrent neural networks.
\newblock In {\em NeurIPS}, 2019.

\bibitem{rubino2021ultra}
A.~Rubino, F.~Rossi, and M.~Imani.
\newblock Ultra-efficient on-device object detection on AI-integrated smart glasses with TinyissimoYOLO.
\newblock In {\em ISLPED}, 2021.

\end{thebibliography}

%═══════════════════════════════════════════════════════════════════════════════
% APPENDIX
%═══════════════════════════════════════════════════════════════════════════════

\appendix

\section{Proof: Closed-Form Solution Derivation}
\label{app:proof}

For the ODE $\frac{d\hvx}{dt} = \frac{\xinf - \hvx}{\tau}$ with constant $\xinf$ and $\tau$, define $\mathbf{y}(t) = \hvx(t) - \xinf$. Then:
\begin{equation}
    \frac{d\mathbf{y}}{dt} = -\frac{\mathbf{y}}{\tau}
\end{equation}
This decouples into $D$ independent scalar ODEs $\frac{dy_i}{dt} = -y_i / \tau$, each with solution $y_i(t) = y_i(0) \cdot e^{-t/\tau}$. Therefore:
\begin{equation}
    \hvx(t + \Delta t) = \xinf + (\hvx(t) - \xinf) \cdot e^{-\Delta t / \tau}
\end{equation}
which is Eq.~\ref{eq:exact}. The computation requires $O(D)$ operations regardless of $\Delta t$. \hfill $\square$

\section{Implementation Details}
\label{app:implementation}

\paragraph{Neuron parameters.} Each neuron stores 6 hypervectors: state $\hvx$, weight $\hvw$, input mask $\hvu$, tau modulator $\mathbf{m}_\tau$, gate weight $\mathbf{g}$, gate bias $\mathbf{b}_g$. Total: $6 \times 16{,}384 \times 4 = 384$KB per neuron. Weight and input mask also have momentum buffers, bringing the total to $8 \times 16{,}384 \times 4 = 512$KB per neuron.

\paragraph{Initialization.} Weight hypervectors are initialized from a deterministic \emph{genesis seed} using BLAKE3 hashing, ensuring reproducibility across runs. Gate biases are initialized to $0.1 \times$ random to start near-neutral interpolation.

\paragraph{State bounds.} States are soft-bounded: if $\|\hvx\| > 5.0$, the state is rescaled to norm 5.0. This prevents numerical explosion while allowing the network to use the full dynamic range.

\paragraph{Network topology.} Default configuration: 3 layers of [4, 8, 4] neurons with layer binding and optional skip connections. Inter-layer communication uses binding with layer-specific random vectors, ensuring distinct representations at each depth.

\section{Additional Benchmark Details}
\label{app:benchmarks}

\begin{table}[h]
\centering
\caption{Extended benchmark results. Results marked $\dagger$ use synthetic data.}
\label{tab:extended}
\begin{tabular}{lcccl}
\toprule
\textbf{Benchmark} & \textbf{Metric} & \textbf{Result} & \textbf{$D$} & \textbf{Notes} \\
\midrule
\multicolumn{5}{l}{\emph{Real-world datasets}} \\
ISOLET Letter Recognition & Accuracy & 91.7\% & 4,096 & 26 classes, 617 features \\
MNIST Digit Classification & Accuracy & 88.5\% & 4,096 & 5 retrain iterations \\
Ethics (Moral Algebra) & Accuracy & 92.9\% & 4,096 & 5 ethical categories \\
Sleep Staging (Sleep-EDF) & Stages & 5/5 & 4,096 & 25.2\% overall (dual-ch.) \\
\midrule
\multicolumn{5}{l}{\emph{Synthetic benchmarks}} \\
Speaker Identification$^\dagger$ & Accuracy & 94.5\% & 4,096 & 10 synth.\ speakers, 13 MFCC \\
EEG Seizure Detection$^\dagger$ & Sensitivity & 100\% & 4,096 & Synth.\ spike-wave patterns \\
Emotion EEG$^\dagger$ & Tests & 6/6 & 4,096 & Synth.\ DEAP-like signals \\
\midrule
\multicolumn{5}{l}{\emph{System benchmarks}} \\
Tokamak Control & Throughput & 59K/s & 16,384 & Real-time, CPU-only \\
FEP Active Inference & Speedup & 7.9$\times$ & 16,384 & vs.\ pymdp baseline \\
\bottomrule
\end{tabular}
\end{table}

\end{document}
